% Originally: content/week-02.tex, the Friday warm-up following \subsection{Sequences and convergence}
% Corsin used this to reopen the topic on the Friday; it needs \cref{def:open_ball}, so it is
% placed here, after the open/closed material, rather than where the week happened to put it.
\section{Unit balls of the three standard metrics}
\label{sec:unit_balls}

\begin{warmup}
What is the shape of $B_1(0) \subseteq \mathbb{R}^2$ with the supremum metric? With the Manhattan
metric?
\end{warmup}

\begin{answer}
Supremum: the axis-aligned \textbf{open square} $(-1,1)^2$. Manhattan: the \textbf{open diamond}
$\{|x_1|+|x_2| < 1\}$. In both cases the strict inequality in \cref{def:open_ball} matters: the
points $(\pm 1,0)$ and $(0,\pm 1)$ sit on the boundary and are \emph{not} in the ball. They are
the corners of the diamond and the edge midpoints of the square, and are drawn hollow below.
\end{answer}


\begin{center}
\begin{tikzpicture}[>=Latex, scale=1.2]
  \begin{scope}[shift={(0,0)}]
    \draw[->] (-1.5,0) -- (1.5,0) node[right] {$x_1$};
    \draw[->] (0,-1.5) -- (0,1.5) node[above] {$x_2$};
    \draw[blue!80!black, thick, dashed, fill=blue!10] (-1,-1) rectangle (1,1);
    % Hollow: (1,0) and (0,1) have d_3 = 1, so they are NOT in the open ball B_1(0).
    \draw[orange!90!black, thick, fill=white] (1,0) circle (2pt);
    \node[orange!90!black, above right, font=\footnotesize] at (1,0) {$(1,0)$};
    \draw[orange!90!black, thick, fill=white] (0,1) circle (2pt);
    \node[orange!90!black, above right, font=\footnotesize] at (0,1) {$(0,1)$};
    \node[below, font=\small, align=center] at (0,-1.6)
      {Supremum metric $d_3$:\\ open square $(-1,1)^2$};
  \end{scope}

  \begin{scope}[shift={(4.4,0)}]
    \draw[->] (-1.5,0) -- (1.5,0) node[right] {$x_1$};
    \draw[->] (0,-1.5) -- (0,1.5) node[above] {$x_2$};
    \draw[blue!80!black, thick, dashed, fill=blue!10] (1,0) -- (0,1) -- (-1,0) -- (0,-1) -- cycle;
    % Same here: d_1((1,0),0) = 1, so the corners are boundary points, not elements.
    \draw[orange!90!black, thick, fill=white] (1,0) circle (2pt);
    \node[orange!90!black, above right, font=\footnotesize] at (1,0) {$(1,0)$};
    \draw[orange!90!black, thick, fill=white] (0,1) circle (2pt);
    \node[orange!90!black, above right, font=\footnotesize] at (0,1) {$(0,1)$};
    \node[below, font=\small, align=center] at (0,-1.6)
      {Manhattan metric $d_1$:\\ open diamond};
  \end{scope}

  \begin{scope}[shift={(8.8,0)}]
    \draw[->] (-1.5,0) -- (1.5,0) node[right] {$x_1$};
    \draw[->] (0,-1.5) -- (0,1.5) node[above] {$x_2$};
    \draw[blue!80!black, thick, dashed, fill=blue!10] (0,0) circle (1);
    \draw[orange!90!black, thick, fill=white] (1,0) circle (2pt);
    \node[orange!90!black, above right, font=\footnotesize] at (1,0) {$(1,0)$};
    \draw[orange!90!black, thick, fill=white] (0,1) circle (2pt);
    \node[orange!90!black, above right, font=\footnotesize] at (0,1) {$(0,1)$};
    \node[below, font=\small, align=center] at (0,-1.6)
      {Standard metric $d_2$:\\ open disk};
  \end{scope}
\end{tikzpicture}
\end{center}

% Generator: Claude Opus 5 (High)
\begin{remark}[The unit ball remembers which metric it came from]
\label{rem:unit_balls_distinguish_metrics}
Three metrics, three unmistakably different unit balls: the square for $d_3$, the disk for
$d_2$, the diamond for $d_1$. This is the quickest way to tell the three apart, and it also makes
\cref{ex:ai_metric_equivalence} visible, since the inclusions
\[B^{d_1}_1(0) \ \subseteq\ B^{d_2}_1(0) \ \subseteq\ B^{d_3}_1(0)\]
are precisely the inequalities $d_3 \leq d_2 \leq d_1$ read backwards. The larger the metric, the
harder it is to be within distance $1$, and hence the smaller the ball. The same family, continued
to all $p \geq 1$, is sketched in \cref{ex:4.3}\textbf{.1}, where the square reappears as the
limiting case $p \to \infty$.

None of this changes which sets are open. All three metrics have the same open sets, as
\cref{ex:ai_metric_equivalence} shows. The balls look different, and the topology does not notice.
\end{remark}



% Originally: content/week-02.tex, end of \section{Compactness}

\begin{aiexercise}[Boundary of the Open Unit Disk]
\label{ex:ai_disk_boundary}
% Generator: Gemini 3.6 Flash (Medium)
Let $D := \{(x,y) \in \mathbb{R}^2 \mid x^2 + y^2 < 1\}$ be the open unit disk in $\mathbb{R}^2$. Show that the boundary of $D$ is the unit circle $\partial D = \{(x,y) \in \mathbb{R}^2 \mid x^2 + y^2 = 1\}$.
\end{aiexercise}



% --- exercises relocated here from the dissolved sheet 2 ---

% Originally: content/exercise-sheets/week-02.tex, problem 2.2
% Source: exercises/Ex2_Analysis2_eng.pdf, p. 1

\begin{exercise}[2.2 --- Multiple choice (closure)]
\label{ex:2.2}
Let $(X,d)$ be a metric space, and $Y_1, Y_2 \subset X$ subsets. Select all the statements below
that are necessarily true.
\begin{enumerate}[label=\textbf{(\alph*)}]
  \item $\overline{Y_1 \cup Y_2} = \overline{Y_1} \cup \overline{Y_2}$
  \item $\overline{Y_1} \cap \overline{Y_2} \subset \overline{Y_1 \cap Y_2}$
  \item $\overline{Y_1 \cap Y_2} \subset \overline{Y_1} \cap \overline{Y_2}$
  \item $\overline{Y_1 \cap Y_2} = \overline{Y_1} \cap \overline{Y_2}$
\end{enumerate}
\exhint[Official hints]{\textbf{(b):} compare with \cref{ex:2.1}.3. \textbf{(d):} play with a set
of three points (a triangle).}
\end{exercise}

% Originally: content/exercise-sheets/week-02.tex, problem 2.3
% Source: exercises/Ex2_Analysis2_eng.pdf, p. 1

\begin{exercise}[2.3 --- Multiple choice (distance from a set)]
\label{ex:2.3}
Let $(X,d)$ be a metric space, and $A \subset X$ a non-empty subset. We define the function
``distance from $A$'' as
\[
  d(\cdot, A) : X \to [0,\infty), \qquad d(x,A) := \inf_{a \in A} d(x,a).
\]
Select all the statements below that are necessarily true.
\begin{enumerate}[label=\textbf{(\alph*)}]
  \item If $A$ is closed and $x \in A^c$, then $d(x,A) > 0$.
  \item The set $M := \{x \in X : d(x,A) \geq 1\}$ is closed in $X$.
  \item For $x, y \in X$, $d(x,A) \leq d(x,y) + d(y,A)$ holds.
  \item If $A^\circ$ is non-empty and $x \in X$, then $d(x,A) = d(x, A^\circ)$.
\end{enumerate}
\exhint[Official hint]{First convince yourself with a drawing that the name of this function is
appropriate; then use the characterization of open/closed sets with sequences.}
\end{exercise}

% Originally: content/exercise-sheets/week-02.tex, problem 2.4
% Source: exercises/Ex2_Analysis2_eng.pdf, p. 1

\begin{exercise}[2.4 --- Boundary, Interior etc.]
\label{ex:2.4}
Determine the interior, closure, and boundary of the following subsets $Y$ of $\mathbb{R}$, for
the standard topology on $\mathbb{R}$. No need to justify the answer.
\begin{center}
\begin{tabular}{ll ll}
(1) & $Y = [0,1]$ & (2) & $Y = \mathbb{Q}$ \\
(3) & $Y = \emptyset$ & (4) & $Y = (0,1)$ \\
(5) & $Y = [-1,1) \setminus \{0\}$ & (6) & $Y = [0,\infty)$ \\
(7) & $Y = \{0\}$ & (8) & $Y = \left\{\tfrac{1}{n} \mid n \in \mathbb{N}\setminus\{0\}\right\}$
\end{tabular}
\end{center}
\end{exercise}
